\documentclass[10pt,a4paper]{amsart}
\usepackage{amsmath,amssymb,amsthm,graphicx}
% Retained minimal preamble for the Stacks Project chapter sites-cohomology.tex (Cohomology on Sites).
% Verbatim excerpt of the chapter preamble: Xy-pic setup (native preamble lines 14-18) and the native
% theorem declarations with the mathematical macros (native preamble lines 154-217). Copied mechanically;
% nothing added, removed or redefined. The wrapper class and the amsmath/amssymb/amsthm/graphicx base
% packages are supplied by the assembler, so the class/externaldocument part of the native preamble is omitted.
% For commutative diagrams we use Xy-pic
\usepackage[all]{xy}

% We use 2cell for 2-commutative diagrams.
\xyoption{2cell}
% Theorem environments.
%
\theoremstyle{plain}
\newtheorem{theorem}[subsection]{Theorem}
\newtheorem{proposition}[subsection]{Proposition}
\newtheorem{lemma}[subsection]{Lemma}

\theoremstyle{definition}
\newtheorem{definition}[subsection]{Definition}
\newtheorem{example}[subsection]{Example}
\newtheorem{exercise}[subsection]{Exercise}
\newtheorem{situation}[subsection]{Situation}

\theoremstyle{remark}
\newtheorem{remark}[subsection]{Remark}
\newtheorem{remarks}[subsection]{Remarks}

\numberwithin{equation}{subsection}

% Macros
%
\def\lim{\mathop{\mathrm{lim}}\nolimits}
\def\colim{\mathop{\mathrm{colim}}\nolimits}
\def\Spec{\mathop{\mathrm{Spec}}}
\def\Hom{\mathop{\mathrm{Hom}}\nolimits}
\def\Ext{\mathop{\mathrm{Ext}}\nolimits}
\def\SheafHom{\mathop{\mathcal{H}\!\mathit{om}}\nolimits}
\def\SheafExt{\mathop{\mathcal{E}\!\mathit{xt}}\nolimits}
\def\Sch{\mathit{Sch}}
\def\Mor{\mathop{\mathrm{Mor}}\nolimits}
\def\Ob{\mathop{\mathrm{Ob}}\nolimits}
\def\Sh{\mathop{\mathit{Sh}}\nolimits}
\def\NL{\mathop{N\!L}\nolimits}
\def\CH{\mathop{\mathrm{CH}}\nolimits}
\def\proetale{{pro\text{-}\acute{e}tale}}
\def\etale{{\acute{e}tale}}
\def\QCoh{\mathit{QCoh}}
\def\Ker{\mathop{\mathrm{Ker}}}
\def\Im{\mathop{\mathrm{Im}}}
\def\Coker{\mathop{\mathrm{Coker}}}
\def\Coim{\mathop{\mathrm{Coim}}}

% Boxtimes
%
\DeclareMathSymbol{\boxtimes}{\mathbin}{AMSa}{"02}

%
% Macros for moduli stacks/spaces
%
\def\QCohstack{\mathcal{QC}\!\mathit{oh}}
\def\Cohstack{\mathcal{C}\!\mathit{oh}}
\def\Spacesstack{\mathcal{S}\!\mathit{paces}}
\def\Quotfunctor{\mathrm{Quot}}
\def\Hilbfunctor{\mathrm{Hilb}}
\def\Curvesstack{\mathcal{C}\!\mathit{urves}}
\def\Polarizedstack{\mathcal{P}\!\mathit{olarized}}
\def\Complexesstack{\mathcal{C}\!\mathit{omplexes}}
% \Pic is the operator that assigns to X its picard group, usage \Pic(X)
% \Picardstack_{X/B} denotes the Picard stack of X over B
% \Picardfunctor_{X/B} denotes the Picard functor of X over B
\def\Pic{\mathop{\mathrm{Pic}}\nolimits}
\def\Picardstack{\mathcal{P}\!\mathit{ic}}
\def\Picardfunctor{\mathrm{Pic}}
\def\Deformationcategory{\mathcal{D}\!\mathit{ef}}
\usepackage{cleveref}
\crefname{originallemma}{original lemma}{original lemmas}
\Crefname{originallemma}{Original lemma}{Original lemmas}
\makeatletter
\newlabel{homology-section-truncations}{{0118}{source}{}{sourcepointer-1}{}}
\newlabel{homology-section-truncations@cref}{{[originallemma][0][]0118}{[1][0][]source}{}{}{}}
\makeatother
\crefname{originallemma}{original lemma}{original lemmas}
\Crefname{originallemma}{Original lemma}{Original lemmas}
\makeatletter
\newlabel{derived-lemma-subcategory-left-resolution}{{05T7}{source}{}{sourcepointer-2}{}}
\newlabel{derived-lemma-subcategory-left-resolution@cref}{{[originallemma][0][]05T7}{[1][0][]source}{}{}{}}
\makeatother
\crefname{originallemma}{original lemma}{original lemmas}
\Crefname{originallemma}{Original lemma}{Original lemmas}
\makeatletter
\newlabel{derived-lemma-colim-hocolim}{{0949}{source}{}{sourcepointer-3}{}}
\newlabel{derived-lemma-colim-hocolim@cref}{{[originallemma][0][]0949}{[1][0][]source}{}{}{}}
\makeatother
\crefname{originallemma}{original lemma}{original lemmas}
\Crefname{originallemma}{Original lemma}{Original lemmas}
\makeatletter
\newlabel{derived-lemma-commutes-with-countable-sums}{{094A}{source}{}{sourcepointer-4}{}}
\newlabel{derived-lemma-commutes-with-countable-sums@cref}{{[originallemma][0][]094A}{[1][0][]source}{}{}{}}
\makeatother
\crefname{originallemma}{original lemma}{original lemmas}
\Crefname{originallemma}{Original lemma}{Original lemmas}
\makeatletter
\newlabel{derived-lemma-projectors-have-images-triangulated}{{05QW}{source}{}{sourcepointer-5}{}}
\newlabel{derived-lemma-projectors-have-images-triangulated@cref}{{[originallemma][0][]05QW}{[1][0][]source}{}{}{}}
\makeatother
\crefname{originallemma}{original lemma}{original lemmas}
\Crefname{originallemma}{Original lemma}{Original lemmas}
\makeatletter
\newlabel{derived-lemma-derived-canonical-delta-functor}{{0152}{source}{}{sourcepointer-6}{}}
\newlabel{derived-lemma-derived-canonical-delta-functor@cref}{{[originallemma][0][]0152}{[1][0][]source}{}{}{}}
\makeatother
\crefname{originallemma}{original lemma}{original lemmas}
\Crefname{originallemma}{Original lemma}{Original lemmas}
\makeatletter
\newlabel{derived-lemma-representable-homological}{{0149}{source}{}{sourcepointer-7}{}}
\newlabel{derived-lemma-representable-homological@cref}{{[originallemma][0][]0149}{[1][0][]source}{}{}{}}
\makeatother
\title{Compact objects in the derived category of a Grothendieck abelian category are summands of finite complexes built from generators}
\author{The Stacks Project Authors}
\date{stacks-project:sites-cohomology.tex:local-snapshot-20261011; GNU FDL 1.2 or later}
\begin{document}
\maketitle
\noindent\textit{Editorial note.} All delivered mathematical text is the author's own, copied from the pinned Stacks Project chapter file sites-cohomology.tex; no statement or proof step was rewritten, reordered or removed.\par\medskip
\noindent\textit{Editorial note.} The proof is the author's complete argument, including the proof of the improvement claim in the last paragraph and the splitting of the idempotent; the author's closing clause is retained verbatim.\par\medskip
\noindent\textit{Editorial note.} Every labelled reference occurring inside the counted spans points to a chapter of the Project that is not part of this card's source tree (Homology and Derived Categories). Each is carried in assembly.external\_references as a non-load-bearing navigation pointer whose source location is the invoking line of the pinned chapter, so no load-bearing mathematics is replaced by a stub.\par\medskip
\noindent\textit{Editorial note.} The checked chapter list for same-name collisions is the set of stacks chapters staged in this repository: chow, cohomology, divisors, etale-cohomology, flat, more-algebra, more-morphisms, morphisms, spaces-more-morphisms, varieties together with this card's chapter sites-cohomology.\par\medskip
\noindent\textit{Editorial note.} The preamble is the chapter's own retained minimal preamble (sources/stacks-sites-cohomology/retained-preamble.tex), a verbatim excerpt of the native preamble: the Xy-pic setup and the native theorem declarations with the mathematical macros.\par\medskip
\noindent\textit{Allowed standard prerequisites (editorial).}\begin{itemize}
\item Compact object of a triangulated category. Same chapter, the definition of a compact object K in D(A) as one for which Hom\_\{D(A)\}(K, -) commutes with direct sums. Supplies the hypothesis and, through the same property, the finite factorisations on which the whole proof rests.
\item Canonical truncations tau\_\{>= n\} and stupid truncations sigma\_\{>= n\}. Homology, the canonical and stupid truncation functors on complexes and their basic properties. Supplies the two filtrations of the representing complex, whose finite stages replace K up to factorisation of the identity.
\item Bounded above resolution by direct sums of generators and the functoriality of the derived category. Derived Categories, the existence of a bounded above resolution whose terms are direct sums of objects of S and the properties of countable direct sums in the derived category. Supplies the representing complex with terms in direct sums of S and the passage from the union of stupid truncations to a factorisation of the identity.
\item Idempotents in a triangulated category and the kernel of an endomorphism. Derived Categories, the splitting of an idempotent e via Ker(e) and Ker(1 - e) and the homological behaviour of represented functors. Supplies the final step, which produces the required direct summand of a finite complex built from S.
\end{itemize}\medskip
\ifdefined\hypertarget\hypertarget{sourcepointer-1}{}\fi
\noindent\textit{Source reference.} Original lemma 0118: sources/stacks-sites-cohomology/pinned-original-source.tex, source lines 13957--13957. Navigation pointer invoked at proof line 13957; the counted span contains the pointer only and the cited result is not reproduced inside it, so the pointer is registered as a non-load-bearing navigation stub in the established form for Stacks material. No bibliography citation is involved.\par\medskip
\ifdefined\hypertarget\hypertarget{sourcepointer-2}{}\fi
\noindent\textit{Source reference.} Original lemma 05T7: sources/stacks-sites-cohomology/pinned-original-source.tex, source lines 13966--13966. Navigation pointer invoked at proof line 13966; the counted span contains the pointer only and the cited result is not reproduced inside it, so the pointer is registered as a non-load-bearing navigation stub in the established form for Stacks material. No bibliography citation is involved.\par\medskip
\ifdefined\hypertarget\hypertarget{sourcepointer-3}{}\fi
\noindent\textit{Source reference.} Original lemma 0949: sources/stacks-sites-cohomology/pinned-original-source.tex, source lines 13973--13973. Navigation pointer invoked at proof line 13973; the counted span contains the pointer only and the cited result is not reproduced inside it, so the pointer is registered as a non-load-bearing navigation stub in the established form for Stacks material. No bibliography citation is involved.\par\medskip
\ifdefined\hypertarget\hypertarget{sourcepointer-4}{}\fi
\noindent\textit{Source reference.} Original lemma 094A: sources/stacks-sites-cohomology/pinned-original-source.tex, source lines 13974--13974. Navigation pointer invoked at proof line 13974; the counted span contains the pointer only and the cited result is not reproduced inside it, so the pointer is registered as a non-load-bearing navigation stub in the established form for Stacks material. No bibliography citation is involved.\par\medskip
\ifdefined\hypertarget\hypertarget{sourcepointer-5}{}\fi
\noindent\textit{Source reference.} Original lemma 05QW: sources/stacks-sites-cohomology/pinned-original-source.tex, source lines 13995--13995. Navigation pointer invoked at proof line 13995; the counted span contains the pointer only and the cited result is not reproduced inside it, so the pointer is registered as a non-load-bearing navigation stub in the established form for Stacks material. No bibliography citation is involved.\par\medskip
\ifdefined\hypertarget\hypertarget{sourcepointer-6}{}\fi
\noindent\textit{Source reference.} Original lemma 0152: sources/stacks-sites-cohomology/pinned-original-source.tex, source lines 14021--14021. Navigation pointer invoked at proof line 14021; the counted span contains the pointer only and the cited result is not reproduced inside it, so the pointer is registered as a non-load-bearing navigation stub in the established form for Stacks material. No bibliography citation is involved.\par\medskip
\ifdefined\hypertarget\hypertarget{sourcepointer-7}{}\fi
\noindent\textit{Source reference.} Original lemma 0149: sources/stacks-sites-cohomology/pinned-original-source.tex, source lines 14023--14023. Navigation pointer invoked at proof line 14023; the counted span contains the pointer only and the cited result is not reproduced inside it, so the pointer is registered as a non-load-bearing navigation stub in the established form for Stacks material. No bibliography citation is involved.\par\medskip
\begin{lemma}
\label{lemma-compact-in-terms-of-generators}
Let $\mathcal{A}$ be a Grothendieck abelian category. Let
$S \subset \Ob(\mathcal{A})$ be a set of objects such that
\begin{enumerate}
\item any object of $\mathcal{A}$ is a quotient of a direct sum
of elements of $S$, and
\item for any $E \in S$ the functor $\Hom_\mathcal{A}(E, -)$
commutes with direct sums.
\end{enumerate}
Then every compact object of $D(\mathcal{A})$ is a direct summand
in $D(\mathcal{A})$ of a finite complex of finite direct sums of
elements of $S$.
\end{lemma}

\begin{proof}
Assume $K \in D(\mathcal{A})$ is a compact object. Represent $K$ by a complex
$K^\bullet$ and consider the map
$$
K^\bullet
\longrightarrow
\bigoplus\nolimits_{n \geq 0} \tau_{\geq n} K^\bullet
$$
where we have used the canonical truncations, see
Homology, Section \ref{homology-section-truncations}.
This makes sense as in each degree the direct sum on the right is finite.
By assumption this map factors through a finite direct sum.
We conclude that $K \to \tau_{\geq n} K$ is zero for at least one $n$,
i.e., $K$ is in $D^{-}(R)$.

\medskip\noindent
We may represent $K$ by a bounded above complex $K^\bullet$ each of whose
terms is a direct sum of objects from $S$, see
Derived Categories, Lemma \ref{derived-lemma-subcategory-left-resolution}.
Note that we have
$$
K^\bullet = \bigcup\nolimits_{n \leq 0} \sigma_{\geq n}K^\bullet
$$
where we have used the stupid truncations, see
Homology, Section \ref{homology-section-truncations}.
Hence by Derived Categories, Lemmas \ref{derived-lemma-colim-hocolim} and
\ref{derived-lemma-commutes-with-countable-sums}
we see that $1 : K^\bullet \to K^\bullet$ factors through
$\sigma_{\geq n}K^\bullet \to K^\bullet$ in $D(R)$.
Thus we see that $1 : K^\bullet \to K^\bullet$ factors as
$$
K^\bullet \xrightarrow{\varphi} L^\bullet \xrightarrow{\psi} K^\bullet
$$
in $D(\mathcal{A})$ for some complex $L^\bullet$ which is bounded and
whose terms are direct sums of elements of $S$. Say $L^i$ is zero for
$i \not \in [a, b]$. Let $c$ be the largest integer $\leq b + 1$ such
that $L^i$ a finite direct sum of elements of $S$ for $i < c$.
Claim: if $c < b + 1$, then we can modify $L^\bullet$ to increase $c$.
By induction this claim will show we have a factorization
of $1_K$ as
$$
K \xrightarrow{\varphi} L \xrightarrow{\psi} K
$$
in $D(\mathcal{A})$ where $L$ can be represented by a finite
complex of finite direct sums of elements of $S$. Note that
$e = \varphi \circ \psi \in \text{End}_{D(\mathcal{A})}(L)$
is an idempotent. By Derived Categories,
Lemma \ref{derived-lemma-projectors-have-images-triangulated}
we see that $L = \Ker(e) \oplus \Ker(1 - e)$.
The map $\varphi : K \to L$ induces an isomorphism with
$\Ker(1 - e)$ in $D(R)$ and we conclude.

\medskip\noindent
Proof of the claim. Write $L^c = \bigoplus_{\lambda \in \Lambda} E_\lambda$.
Since $L^{c - 1}$ is a finite direct sum of elements of $S$
we can by assumption (2) find a finite subset
$\Lambda' \subset \Lambda$ such that $L^{c - 1} \to L^c$ factors
through $\bigoplus_{\lambda \in \Lambda'} E_\lambda \subset L^c$.
Consider the map of complexes
$$
\pi :
L^\bullet
\longrightarrow
(\bigoplus\nolimits_{\lambda \in \Lambda \setminus \Lambda'} E_\lambda)[-i]
$$
given by the projection onto the factors corresponding to
$\Lambda \setminus \Lambda'$ in degree $i$.
By our assumption on $K$ we see that, after possibly replacing $\Lambda'$ by
a larger finite subset, we may assume that $\pi \circ \varphi = 0$
in $D(\mathcal{A})$.
Let $(L')^\bullet \subset L^\bullet$ be the kernel of $\pi$.
Since $\pi$ is surjective we get a short exact sequence of complexes,
which gives a distinguished triangle in $D(\mathcal{A})$ (see
Derived Categories, Lemma \ref{derived-lemma-derived-canonical-delta-functor}).
Since $\Hom_{D(\mathcal{A})}(K, -)$ is homological (see
Derived Categories, Lemma \ref{derived-lemma-representable-homological})
and $\pi \circ \varphi = 0$, we can find a morphism
$\varphi' : K^\bullet \to (L')^\bullet$ in $D(\mathcal{A})$ whose
composition with $(L')^\bullet \to L^\bullet$ gives $\varphi$.
Setting $\psi'$ equal to the composition of $\psi$ with
$(L')^\bullet \to L^\bullet$ we obtain a new factorization.
Since $(L')^\bullet$ agrees with $L^\bullet$ except in degree $c$
and since $(L')^c = \bigoplus_{\lambda \in \Lambda'} E_\lambda$ the
claim is proved.
\end{proof}

\end{document}
